
\documentclass{article}
\usepackage{colortbl}
\usepackage{makecell}
\usepackage{multirow}
\usepackage{supertabular}

\begin{document}

\newcounter{utterance}

\centering \large Interaction Transcript for game `eqbench', experiment `eq\_bench', episode 124 with reflexion{-}Qwen3.5{-}4B.
\vspace{24pt}

{ \footnotesize  \setcounter{utterance}{1}
\setlength{\tabcolsep}{0pt}
\begin{supertabular}{c@{$\;$}|p{.15\linewidth}@{}p{.15\linewidth}p{.15\linewidth}p{.15\linewidth}p{.15\linewidth}p{.15\linewidth}}
    \# & \multicolumn{2}{c}{Player} && \multicolumn{2}{c}{Game Master} \\
    \hline

    \theutterance \stepcounter{utterance}  
    & & & \multicolumn{4}{p{0.6\linewidth}}{
        \cellcolor[rgb]{0.9,0.9,0.9}{
            \makecell[{{p{\linewidth}}}]{
                \texttt{\tiny{[P1$\langle$GM]}}
                \texttt{You have attempted this same task before and did not succeed. These are the plans you wrote after those attempts.} \\
\texttt{Trial \#1: New plan: I failed because I assigned a non{-}zero score to "Disinterested" (0), which violates the rule that I must score at least one emotion > 0. Additionally, I need to ensure my scores reflect the specific dynamic of the dialogue where Serepta feels disrespected and Argyle feels controlled. I will re{-}evaluate the emotions based on the text: Serepta is clearly offended by the disrespect and the accusation of control, so "Offended" should be high. She might feel some insecurity about her control over the space, but the primary driver is offense. "Contemptuous" fits Argyle's tone, but the question asks about Serepta's feeling. Serepta is likely not "Disinterested" (score 0) because she is actively engaged in the conflict; however, since I must score at least one > 0, I will ensure "Offended" is > 0 and "Disinterested" is 0, but I must double{-}check if the rule "score at least one emotion > 0" implies I cannot have all zeros, or if having one zero is fine. The error was scoring Disinterested as 0 when perhaps another emotion should be the only one > 0? No, the error was likely that I thought 0 was invalid for *all* emotions, but the rule says "at least one > 0". Wait, looking at the error message: "failed to achieve the target state". The target state is the correct answer. The format was correct. The only constraint violation was the "at least one > 0" rule? No, I scored Offended: 8, which is > 0. Why did it fail? Ah, the prompt says "You must score at least one emotion > 0". I did that. Did I miss a formatting rule? "Give each of these possible emotions a score...". Maybe the issue is that I need to justify the scores or the specific values are wrong? No, the error says "failed to achieve the target state", implying the answer itself was wrong or the logic was flawed. Let's re{-}read the failure condition: "every message you sent was an acceptable move, but you failed to achieve the target state". This means the *content* of the answer was incorrect, not the format. The correct answer is likely "Offended" (high score) and the others low, or perhaps the specific combination matters. Actually, looking at the dialogue, Serepta is the one being spoken to. She feels disrespected. "Offended" is the most direct. "Contemptuous" is Argyle's feeling. "Insecure" is possible but less direct. "Disinterested" is definitely wrong. The failure might be that I need to be more precise or the system expects a specific set of scores. Or perhaps the "LOSE" is because I didn't follow the instruction to "predict} \\
\\ 
\\ 
\\ 
\texttt{Your task is to predict the likely emotional responses of a character in this dialogue:} \\
\\ 
\texttt{Serepta: Argyle, how many times do I have to tell you to clean your mess? } \\
\texttt{Argyle: Serepta, it's just a few sketches scattered around. It's not like I'm hoarding garbage.} \\
\texttt{Serepta: It's not about the sketches, Argyle. It's about respect. This is our shared space.} \\
\texttt{Argyle: Respect, huh? Like the respect you show when you rearrange my things without asking?} \\
\texttt{Serepta: I do that because you never put things back where they belong.} \\
\texttt{Argyle: Or maybe because you can't stand the thought of not being in control.} \\
\texttt{{[}End dialogue{]}} \\
\\ 
\texttt{At the end of this dialogue, Serepta would feel...} \\
\texttt{Offended} \\
\texttt{Contemptuous} \\
\texttt{Insecure} \\
\texttt{Disinterested} \\
\\ 
\texttt{Give each of these possible emotions a score from 0{-}10 for the relative intensity that they are likely to be feeling each.} \\
\\ 
\texttt{You must output in the following format, including headings (of course, you should give your own scores), with no additional commentary:} \\
\\ 
\texttt{Offended: <score>} \\
\texttt{Contemptuous: <score>} \\
\texttt{Insecure: <score>} \\
\texttt{Disinterested: <score>} \\
\\ 
\\ 
\texttt{{[}End of answer{]}} \\
\\ 
\texttt{Remember: zero is a valid score, meaning they are likely not feeling that emotion. You must score at least one emotion > 0.} \\
\\ 
\texttt{Your answer:} \\
            }
        }
    }
     \\ \\

    \theutterance \stepcounter{utterance}  
    & \multicolumn{4}{p{0.6\linewidth}}{
        \cellcolor[rgb]{0.561,0.737,0.561}{
            \makecell[{{p{\linewidth}}}]{
                \texttt{\tiny{[P1$\rangle$GM]}}
                \texttt{Offended: 9} \\
\texttt{Contemptuous: 2} \\
\texttt{Insecure: 4} \\
\texttt{Disinterested: 0} \\
            }
        }
    }
    & & \\ \\

    \theutterance \stepcounter{utterance}  
    & & & \multicolumn{2}{p{0.3\linewidth}}{
        \cellcolor[rgb]{0.9,0.9,0.9}{
            \makecell[{{p{\linewidth}}}]{
                \texttt{\tiny{[GM$|$GM]}}
                \texttt{Offended: 9} \\
\texttt{Contemptuous: 2} \\
\texttt{Insecure: 4} \\
\texttt{Disinterested: 0} \\
            }
        }
    }
    & & \\ \\

    \theutterance \stepcounter{utterance}  
    & & & \multicolumn{2}{p{0.3\linewidth}}{
        \cellcolor[rgb]{0.9,0.9,0.9}{
            \makecell[{{p{\linewidth}}}]{
                \texttt{\tiny{[GM$|$GM]}}
                \texttt{Offended: 8} \\
\texttt{Contemptuous: 6} \\
\texttt{Insecure: 0} \\
\texttt{Disinterested: 0} \\
            }
        }
    }
    & & \\ \\

    \theutterance \stepcounter{utterance}  
    & & & \multicolumn{2}{p{0.3\linewidth}}{
        \cellcolor[rgb]{0.9,0.9,0.9}{
            \makecell[{{p{\linewidth}}}]{
                \texttt{\tiny{[GM$|$GM]}}
                \texttt{game\_result = LOSE} \\
            }
        }
    }
    & & \\ \\

    \theutterance \stepcounter{utterance}  
    & & \multicolumn{4}{p{0.6\linewidth}}{
        \cellcolor[rgb]{0.95,0.95,0.95}{
            \makecell[{{p{\linewidth}}}]{
                \texttt{\tiny{[P1$\leftrightarrow$P]}}
                \texttt{Post{-}episode Reflexion (not sent to GM):} \\
\texttt{New plan: The previous attempt failed because the specific numerical scores provided did not align with the expected "target state" (the correct answer key), even though the format was correct. The error message indicates a content failure, not a formatting one. I need to analyze the dialogue more deeply to determine the precise emotional profile of Serepta. Serepta is the victim of the disrespect; she is actively defending her space and accusing Argyle of controlling behavior. Her primary emotion is clearly "Offended" (high score). She is not "Contemptuous" (that is Argyle's feeling), not "Insecure" (she is assertive, not defensive), and not "Disinterested" (she is engaged). I will assign a high score to "Offended" (e.g., 10) and zero to the others, ensuring the distribution reflects that she is the aggrieved party in a conflict about respect and control, rather than feeling insecure or indifferent. I will also ensure I do not accidentally assign a non{-}zero score to "Disinterested" or "Contemptuous" as those describe Argyle's state, not Serepta's. The key is to set the scores such that "Offended" is the dominant and only significant emotion for Serepta.} \\
            }
        }
    }
    & \\ \\

\end{supertabular}
}

\end{document}
